\section{Purpose}

\textbf{[STIPULATION]} In FFGFT $\alpha$ follows from the geometry (A130). This
document treats the complementary question: how does one set $\alpha=1$ \emph{as
a computational tool} in a way that is fully reversible and makes the coupling
structure visible? The answer is the Heaviside--Lorentz stipulation $e^2=4\pi\alpha$
combined with $\alpha=1$ -- not a unit choice but a \emph{stipulation} with
preserved trace.

\section{Stipulation: Heaviside--Lorentz with $\alpha = 1$}

\textbf{[STIPULATION]} In natural units ($\hbar=c=1$) and the Heaviside--Lorentz
system, the \emph{defining equation} is
\[
  e^2 = 4\pi\alpha.
\]
The fine-structure constant is set to $\alpha=1$. This is not a \emph{norm} (unit
choice) -- a dimensionless number is invariant under rescaling of base units --
but a \emph{stipulation}: the coupling is adopted as a computational unit.

\textbf{[B]} The charge carries the $4\pi$; what is set to $1$ is $\alpha$, not
$e$. From $e^2=4\pi\alpha$ with $\alpha=1$ follows the value
$e=\sqrt{4\pi}\approx3{.}545$ -- not $e=1$. The relation $e^2=4\pi\alpha$
remains as a defining equation and cannot be dissolved by the stipulation.

\section{Back-Computation as Substitution}

\textbf{[B]} Because $e^2=4\pi\alpha$ remains as a relation, the coupling power
is automatically carried along. Back-computation is pure substitution
\[
  \alpha \;\longrightarrow\; \frac{1}{137{.}036},
\]
applied to the $\alpha$-order of the process:
\begin{itemize}[leftmargin=2em,itemsep=2pt,topsep=2pt]
  \item per vertex a factor $e=\sqrt{4\pi\alpha}$; compared to the $\alpha{=}1$
    value $\sqrt{4\pi}$, that is a $\sqrt\alpha$ back,
  \item amplitude with $V$ vertices $\propto(4\pi\alpha)^{V/2}
    \;\Rightarrow\;\times\alpha^{V/2}$,
  \item rate/cross section $\propto(4\pi\alpha)^V
    \;\Rightarrow\;\times\alpha^V$,
  \item compact: process of rate order $\alpha^n
    \;\Rightarrow\;\times\alpha^n$.
\end{itemize}

\textbf{[B]} The $4\pi$ is not double-counted: it is already in $e^2$ and in the
phase-space normalisation. Only $\alpha^n$ is reset, not $(4\pi\alpha)^n$.

\section{Conversion Examples}

\textbf{[K]} All quantities in $\hbar=c=1$, Heaviside--Lorentz.
Back-factor $=$ the $\alpha$-order of the expression.

\begin{center}
\renewcommand{\arraystretch}{1.35}
{\small
\begin{tabular}{lllc}
\toprule
Quantity & Physical form & With $\alpha{=}1$ & Back-factor \\
\midrule
Coulomb potential      & $V = \alpha/r$                          & $1/r$            & $\times\alpha$      \\
Class.\ electron radius & $r_e = \alpha/m_e$                     & $1/m_e$          & $\times\alpha$      \\
Compton wavelength     & $\lambda_C = 1/m_e$                     & $1/m_e$          & $\times 1$          \\
Bohr radius            & $a_0 = 1/(\alpha m_e)$                  & $1/m_e$          & $\times\alpha^{-1}$ \\
Rydberg / binding      & $E = \tfrac12\alpha^2 m_e$              & $\tfrac12 m_e$   & $\times\alpha^2$    \\
Thomson cross section  & $\sigma_T = \tfrac{8\pi}{3}\alpha^2/m_e^2$ & $\tfrac{8\pi}{3}/m_e^2$ & $\times\alpha^2$ \\
\bottomrule
\end{tabular}
}
\renewcommand{\arraystretch}{1.0}
\end{center}

\section{The 137-Ladder: What $\alpha=1$ Makes Visible}

\textbf{[B]} The three atomic length scales stand in the ratio
\[
  r_e : \lambda_C : a_0 \;=\; \alpha : 1 : \alpha^{-1}.
\]
At $\alpha=1$ they all collapse to $1/m_e$: they are not three independent
scales but \emph{one} scale with three $\alpha$-powers. $a_0$ is exactly
$137\times$ larger than $\lambda_C$, and $\lambda_C$ exactly $137\times$ larger
than $r_e$.

\textbf{[B]} With $\alpha=1/137$ written out everywhere, this common origin
remains hidden -- three different awkward numbers appear and the $\alpha$-power
as ordering principle is not visible. The $\alpha=1$ stipulation reveals it.

\section{Why This Stipulation -- Comparison with Alternatives}

\textbf{[S]} \textbf{Versus Gaussian ($e^2=\alpha$).} Also works, but places the
$4\pi$ in the field equations rather than in Coulomb's law. Heaviside--Lorentz
places the $4\pi$ where they geometrically belong (sphere surface, Coulomb) and
leaves the Maxwell equations $4\pi$-free. $\alpha$ is then the pure coupling, not
a mixture of coupling and geometry.

\textbf{[S]} \textbf{Versus $c=\hbar=1$ unit choice.} The setting $c=\hbar=1$
destroys the powers -- the trace is gone, back-computation runs only via
dimensional analysis (A266). The defining equation $e^2=4\pi\alpha$, by contrast,
\emph{retains the trace}: the coupling power is carried along because the relation
stands. Stipulation with preserved trace, not norm with erased trace -- hence
constructively reversible, not merely justifiably reconstructable.

\section{FFGFT Connection: Back-Computing $\alpha$ Means Inserting $\xi$}

\textbf{[B]} $\alpha$ is in the FFGFT corpus not a free input but follows from
the geometry (A130):
\[
  \alpha^{-1} = 137{.}036 \quad\text{from}\quad \xi = \frac{4}{30000},
  \quad D_f = 3-\xi, \quad E_0=\sqrt{m_em_\mu}.
\]
The stipulation $\alpha=1$ isolates precisely the point at which the $\xi$-geometry
enters. Every reset $\alpha^n$ is a location at which $D_f=3-\xi$ becomes visible.

\textbf{[B]} Within the corpus, ``back-computing $\alpha$'' therefore does not
mean ``inserting $1/137$'' but ``inserting the $\xi$-structure'' -- and the
$\alpha{=}1$ system makes countable at how many places and in which power that
occurs.

\section{Symbol Glossary}

Symbols used throughout the entire A-series are explained in A010.
Specific to this document:

\begin{center}
\renewcommand{\arraystretch}{1.3}
{\small\begin{tabular}{lp{9cm}}
\toprule
Symbol & Meaning \\
\midrule
$e^2 = 4\pi\alpha$ & Heaviside--Lorentz defining equation \\
$e = \sqrt{4\pi}\approx3{.}545$ & Elementary charge at $\alpha=1$ \\
$r_e = \alpha/m_e$ & Classical electron radius \\
$\lambda_C = 1/m_e$ & Compton wavelength \\
$a_0 = 1/(\alpha m_e)$ & Bohr radius \\
$V$ & Number of Feynman vertices in the process \\
$\alpha^n$ & $\alpha$-order / back-factor of the process \\
\bottomrule
\end{tabular}}
\renewcommand{\arraystretch}{1.0}
\end{center}

\section{Verification Script}

\begin{itemize}[leftmargin=2em,itemsep=2pt,topsep=2pt]
  \item Numerical verification of all table rows with $\alpha=1/137{.}036$
    and $m_e=0{.}511\,\text{MeV}$; ratio $r_e:\lambda_C:a_0=\alpha:1:\alpha^{-1}$:
    \texttt{python/A\_Serie\_Skripte/a267\_alpha\_eins.py}
\end{itemize}

\section{Open Questions}

\textbf{The stipulation $\alpha=1$ is fully worked out only for electromagnetic
processes.} For weak and strong coupling ($\alpha_W$, $\alpha_s$) an analogous
table would be possible but has not been worked out.

\textbf{The connection $\alpha^n$ as locus of $D_f=3-\xi$} is plausibly shown
for simple QED processes but not for multi-loop or non-abelian processes.

\section{Supersedes}

This document subsumes: Dok.~011 (individual formulas $r_e$, $\lambda_C$, $a_0$,
dimensional analysis $\alpha=r_e/\lambda_C$, SI numerical values of atomic
scales), Dok.~015 (Heaviside--Lorentz line $e^2/(4\pi)=1$ in the unit
systematics). The coherent conversion table with explicit $\alpha^n$ back-factors
and the 137-ladder $r_e:\lambda_C:a_0=\alpha:1:\alpha^{-1}$ as a structural
statement are not present as a unit in the legacy corpus and are newly
systematised here. Related: A130 (geometric derivation $\alpha$ from $\xi$),
A266 (SI back-computation).

\section{Use of AI Tools}

This document was created with the assistance of AI tools. Responsibility
for content and errors rests with the author. Full wording of the rules in A010.
